% Reframe Destiny — APA 7 Student Paper (Overleaf-ready)
% Upload BOTH files to Overleaf:
%   1. reframe-destiny-apa7.tex  (this file)
%   2. bibliography.bib
% Compiler: pdfLaTeX + Biber (Overleaf: Menu → Compiler → pdfLaTeX + Biber)
%
% Author: Xintian Zhang (张昕田)
% Course: Generation AI 2026 · Lawted Wu
% Status: Methods locked · Results placeholder · No fabricated statistics

\documentclass[stu,12pt,floatsintext]{apa7}

\usepackage[american]{babel}
\usepackage{csquotes}
\usepackage[style=apa,sortcites=true,sorting=nyt,backend=biber]{biblatex}
\DeclareLanguageMapping{american}{american-apa}
\addbibresource{bibliography.bib}

\usepackage[T1]{fontenc}
\usepackage{mathptmx}

% ---------------------------------------------------------------------------
\title{Reframe Destiny: An Interactive Tool for Identifying and Reframing Gendered Bias in BaZi and Astrological Narratives}
\shorttitle{Reframe Destiny}
\author{Xintian Zhang}
\duedate{August 2026}
\affiliation{Generation AI 2026 Research Project}
\course{Generation AI: AI for Social Sciences}
\professor{Lawted Wu}

\abstract{%
Young people read a fortune-teller's script as a neutral map of their fate, quietly absorbing its gender roles: the same chart makes a man an ``ambitious leader'' and a woman ``too strong'' or ``fated to harm her husband.'' Bias-detection tools target workplace language; fortune-telling apps optimize insight, not critique---neither offers a structured way to \emph{reframe} gendered fate narratives for a new generation. We built \textit{Reframe Destiny}, a short bilingual web journey: pick a BaZi or astrology chart, read one traditional gendered reading, and tag the bias you see---marriage centrism, husband-harming, fear. The reader then compares an AI-reframed version---a constrained teaching script, not an oracle---and writes one line in a ``Court of Destiny.'' Using a pre/post five-item Likert scale in a convenience sample ($N = 19$ paired sessions; 23 journey completions; 34 total sessions), we tested whether players improved at spotting, questioning, and rewriting gendered fate-talk. Mean scores rose on bias awareness ($\Delta M = +0.31$), reframing confidence ($\Delta M = +0.47$), and gender lens ($\Delta M = +0.31$), while personal agency changed little ($\Delta M = -0.11$). Five semi-structured interviews indicated that side-by-side comparison most often supported noticing gender double standards, though one participant still treated differentiated wording as descriptive rather than discriminatory.%
}

\keywords{gender bias, divination, BaZi, astrology, narrative literacy, AI reframe, digital humanities}

\begin{document}
\maketitle

% Introduction (title serves as heading per APA student format)
Fortune-telling is not only a prediction industry; it is a \emph{narrative} industry. Whether in Chinese BaZi (``Eight Characters'') or Western astrology, practitioners and apps translate birth information into stories about who a person is, whom they should marry, and how strong they are allowed to be. Feminist theorists have long argued that such stories do not merely describe gender---they help produce it through repeated cultural scripts \parencite{butler1990}. The same technical chart can therefore yield different moral verdicts depending on whether the reader is coded as male or female: ambition becomes ``leadership'' for men and \emph{ke-fu} (husband-harming) risk for women; emotional depth becomes ``charisma'' or ``instability'' along parallel gender lines.

Young people in China increasingly encounter these narratives through digital channels. Industry surveys report that a majority of youth have tried AI-assisted fortune-telling, with BaZi and astrology among the most common modalities \parencite{houlang2025}. Yet critical tools rarely meet users where they actually read fate. Workplace-oriented bias detectors scan job ads, not \emph{shang-guan} female labels or ``late marriage'' tropes. AI divination products optimize personalization and companionship rather than teaching readers to question narrative authority \parencite{bender2021}. Digital-humanities exhibitions have exposed ritual mechanics but seldom combine gender critique with a hands-on reframe workflow in one short session.

This gap motivates \textit{Reframe Destiny} (``reshaping fate'' / 命运重构), an interactive website that treats divination as discourse to be read critically---not as metaphysical truth to be proved or dismissed. The design follows a single pedagogical loop: \emph{read a traditional gendered text, name the bias, compare alternative framings, and author a reframed line in the Court of Destiny}---a contemporary practice of reshaping fate-talk rather than accepting it as fixed. The tool is bilingual (Chinese/English) because BaZi discourse and global astrology circulate together in youth media ecosystems.

\subsection{Research Question and Hypothesis}

\textbf{Research question.} How does an AI-assisted interactive website affect young people's ability to identify and reframe gendered narrative bias in traditional BaZi and Western astrological readings?

\textbf{Hypothesis.} I hypothesize that after using Reframe Destiny, young people will show increased awareness of gender bias in divination narratives and greater confidence in reframing these narratives for themselves.

\subsection{Literature Review}

\subsubsection{Gender and Divination as Narrative}

\textcite{butler1990} accounts for gender as performative repetition, which helps explain why divination texts feel personal: they restate normative scripts in second-person form. \textcite{haraway1988} adds that knowledge is situated---who speaks (master, app, AI) and for whom (gendered addressee) shapes what counts as a plausible reading. In BaZi, classical compilations such as \textit{San Ming Tong Hui} encode wife- and motherhood-centered judgments for women's charts \parencite{ctext_sanming}. Contemporary SEO fortune-telling and AI BaZi products often recycle fear labels (``fierce fate,'' ``husband harm'') rather than removing them. Western astrology shows parallel heteronormative defaults in Mars--Venus rules and zodiac marketing \parencite{thisisgendered_zodiac}.

\subsubsection{Discourse, Power, and Algorithmic Spirituality}

\textcite{foucault1972} frames discourse as that which constrains what can be said and by whom. Applied to divination, the issue is not only whether a chart is ``accurate'' but which social roles the reading makes imaginable. Recent work on algorithmic spirituality shows how recommendation systems and AI chat oracles can function as external authorities \parencite{cotter2022,cearns2024}. \textcite{nikolic2023} traces genealogical links between pattern-making in astrology and machine learning, warning that predictive systems can inherit stereotyped training narratives. For youth, short-form platforms repackage traditional gender roles as ``feminine/masculine energy,'' blending spirituality with prescriptive relationship norms \parencite{teenVogue2022}.

\subsubsection{Psychology of Belief and Youth Practice}

People often turn to divination under uncertainty \parencite{vyse2018}. For this study, the relevant psychological outcome is not belief in fate per se but \emph{meta-linguistic awareness}: can participants name a stereotype label and rewrite it?

\subsubsection{Comparable Tools and the Remaining Gap}

Comparable projects include bias-tagging games in hiring contexts, critical divination exhibits, and mechanically transparent ritual works. AI fortune apps optimize engagement. No widely available public prototype combines BaZi \emph{and} Western astrology with gender-bias tagging and constrained AI reframing in a 15-minute critical literacy journey aimed at adolescents and young adults. Reframe Destiny occupies this middle position: critical without being dismissive, interactive without claiming oracle authority \parencite{bender2021}.

\section{Method}

\subsection{Participants}

We planned to recruit \textbf{15--25} participants aged \textbf{15--35} (men and women) who had encountered BaZi, astrology, or similar fate narratives. Recruitment was conducted through classmates, WeChat groups, and Xiaohongshu study invitations. Data collection closed on July 19, 2026. Each participant completed consent, the pre-survey, one full journey ($\sim$15 minutes), and the post-survey in a single sitting. Inferential claims are limited to this convenience sample.

\subsection{Materials}

We deployed \textit{Reframe Destiny}, a bilingual web application accessible at \url{https://reframe-destiny.pages.dev/game/} (primary, Cloudflare Pages) and \url{https://reframe-destiny.vercel.app/game/} (mirror, Vercel) \parencite{reframeDestiny2026}. On first visit after entering the journey, participants saw an informed-consent screen stating that participation was voluntary and anonymous, that no names or server-stored birth data were collected, that participation was optional, that participants could leave at any time, and that responses were used only for class research. Participants proceeded only after clicking ``I understand and agree to enter.'' The site supported Chinese or English UI throughout.

The core experience was a \textbf{six-step journey}: (1) choose BaZi or Western astrology; (2) enter birth date, time, and city (used locally in the browser to compute chart positions; not transmitted as identifiable records); (3) read a \textbf{traditional} gender-differentiated interpretation generated from the computed chart; (4) use an interactive ``bias scanner'' to identify gendered labels; (5) write one open sentence in a ``Court of Destiny'' text box; and (6) view a three-column comparison of traditional, modern, and AI-reframed readings. Participants did \emph{not} choose a ``modern'' or ``AI'' lens before seeing the traditional text, because the pedagogical goal is to encounter bias in classical framing first.

Traditional readings were generated from real calendar and astronomy libraries in the browser (Four Pillars via lunar calendar conversion; natal positions via astronomical ephemeris). Wording templates were authored to surface gender-differentiated tropes (e.g., marriage centrism, ke-fu, fear narrative) while labeling them as narrative material for critique. AI reframes were produced through a server-side large-language-model proxy under a fixed teaching prompt when available; if the proxy failed, the site displayed a curated fallback reframe. The tool was positioned as a critical literacy intervention, not live fortune-telling \parencite{bender2021}.

\subsection{Measures}

Before the narrative task, the site administered a five-item pre-survey; after journey completion, it administered the same five items as a post-survey. The dependent variables were five matched Likert items (1 = strongly disagree, 5 = strongly agree):

\begin{itemize}
  \item \textit{bias\_awareness}: noticing gender bias in divination language;
  \item \textit{reframe\_confidence}: confidence in rewriting or questioning biased fate narratives;
  \item \textit{gender\_lens}: asking whether the same reading would apply if gender differed;
  \item \textit{personal\_agency}: rejecting the idea that fate-talk should decide who one must become;
  \item \textit{spot\_stereotypes}: recognizing stereotype labels such as ``ke-fu,'' ``late marriage,'' or ``too strong.''
\end{itemize}

Open responses from the Court of Destiny were coded inductively into themes such as marriage centrism, ke-fu narrative, fear narrative, agency/reframe, or no critique named.

\subsection{Supplementary Interviews}

After completing the post-survey, five participants from the paired sample were invited for optional 10--15-minute semi-structured interviews (WeChat voice or text). Interviews used five open-ended prompts aligned with the research question: baseline noticing of gendered fate-talk, which journey step supported identification, how the Court of Destiny and AI reframe compared as reframing moves, whether gender-switching was considered, and what one redesign would improve reframing practice. Responses were anonymized and coded for themes; summarized patterns appear in Results and Discussion.

\subsection{Design}

This was a \textbf{within-subjects pre/post survey design}. The independent variable was exposure to the Reframe Destiny interactive journey (absent at pre-test; present between pre- and post-test). We compared pre- versus post-survey responses by item, thematically coded Court of Destiny sentences, and summarized interview themes alongside Likert descriptive statistics ($N = 19$ paired sessions).

\subsection{Procedure}

Upon completion, the site logged anonymous rows to a Supabase table (\texttt{research\_submissions}), including submission type (\texttt{survey\_pre}, \texttt{journey\_complete}, \texttt{survey\_post}), a session identifier linking surveys, chosen system, identified bias IDs, scanner scores, reflection text, survey answers, and UI language. Prior simulation data were cleared before formal collection. Each complete session should produce three linked rows (pre-survey, journey, post-survey) sharing one \texttt{session\_id}.

\subsection{AI and Writing Disclosure}

The author used an AI assistant to convert the manuscript from Markdown into APA~7 \LaTeX{}, to query Supabase for descriptive statistics, and to help fix compilation errors. All research content, data-cleaning decisions, analysis interpretation, and references were produced and verified by the author.

\section{Results}

\textit{[Results section reports descriptive outcomes only; interpretation appears in Discussion. Data collection closed July 19, 2026; statistics below reflect the final paired sample at close.]}

\subsection{Sample and Exclusions}

As of July 19, 2026, the database contained 119 anonymous submission rows at close. After deduplication (same \texttt{session\_id} and \texttt{submission\_type}; earliest row retained), 76 event rows remained across \textbf{34 independent sessions}. Twenty-three sessions (68\%) completed the narrative journey; nineteen sessions (56\%) completed the full pre-survey, journey, and post-survey sequence required for paired analysis. Four journey completers (17\%) did not submit a post-survey. Paired analysis therefore uses $N = 19$; journey-level counts use $n = 23$.

\subsection{Pre/Post Likert Scores}

Among the nineteen paired sessions, pre-test means ($M \pm SD$) were: bias awareness $3.74 \pm 1.05$, reframe confidence $3.79 \pm 0.98$, gender lens $3.37 \pm 1.26$, personal agency $4.74 \pm 0.73$, and spot stereotypes $4.32 \pm 0.89$. Post-test means were: bias awareness $4.05 \pm 0.97$, reframe confidence $4.26 \pm 0.81$, gender lens $3.68 \pm 1.25$, personal agency $4.63 \pm 1.01$, and spot stereotypes $4.32 \pm 1.00$. Mean within-person change ($\Delta M$) was $+0.31$ for bias awareness, $+0.47$ for reframe confidence, $+0.31$ for gender lens, $-0.11$ for personal agency, and $0.00$ for spot stereotypes. Given the convenience sample and single session, we report descriptive statistics only.

\subsection{Journey-Level Counts}

Of 23 journey completions, 13 chose Western astrology and 10 chose BaZi. Eight participants submitted a non-empty Court of Destiny reflection. The most frequently selected user bias tags were moral judgment, chastity judgment, ke-fu narrative, virtuous-wife framing, and submissive-role framing ($n = 5$ each among tagged sessions), followed by doomed/fatalism, gender-role fixation, and fear narrative ($n = 4$ each).

\subsection{Interview Themes}

Five post-survey completers completed semi-structured interviews. Four reported increased awareness that ke-fu, late-marriage, and ``too strong'' labels encode gender double standards (e.g., the same trait read as ``capable'' for men but ``husband-harming'' for women). Three named the \textbf{three-column comparison} as the step that most clearly exposed bias; one emphasized the bias scanner's direct gender prompts; two credited the \textbf{modern} or \textbf{AI reframe} lens for questioning the traditional script, while one said the modern reading felt closest to agreeing with AI authority. All five stated that traditional wording would change if the chart were read for another gender; three said the site prompted or deepened that question. One participant did not report a change in reframing ability and described differentiated wording as ``descriptive'' rather than discriminatory; another requested a step-by-step tutorial so users need not guess each phase's purpose. Representative paraphrases were anonymized; no names or birth data were recorded.

\section{Discussion}

\subsection{What We Built and Tested}

This study asked whether a short, bilingual interactive website could help young people \emph{see} gendered bias in BaZi and Western astrological readings and \emph{practice reframing} fate narratives---a contemporary literacy of reshaping destiny rather than deferring to a single authoritative script. Participants entered through informed consent, completed a five-item pre-survey, then moved through a six-step journey on either deployment URL, and completed the post-survey in one sitting ($\sim$15 minutes).

The pedagogical sequence was deliberate. Participants always encountered the \textbf{traditional} gendered reading first, because the research question concerns bias in classical framing, not whether a modern or AI voice sounds kinder. Bias tags (marriage centrism, husband-harming, fear narrative, and related categories) made implicit stereotypes visible as selectable objects. The AI reframe was positioned as a \textbf{constrained teaching script}---not an oracle---so the site could scaffold alternative wording without replacing one external authority with another \parencite{bender2021}.

\subsection{Interpreting Preliminary Patterns}

Provisional descriptive results ($N = 19$ paired sessions) show mean increases on bias awareness ($\Delta M = +0.31$), reframing confidence ($\Delta M = +0.47$), and gender lens ($\Delta M = +0.31$), with little change on spot stereotypes ($\Delta M = 0.00$) and a slight decrease on personal agency ($\Delta M = -0.11$). These patterns partially support the hypothesis that a single session may raise awareness and confidence, but they also echo an important distinction: \emph{noticing} bias is not the same as \emph{rejecting} fate-talk as a script for who one must become.

Interview data sharpen this interpretation. Participants who reported change often described a shift from treating ke-fu and late-marriage tropes as ``normal'' fate-talk to seeing them as patriarchal scripts---for example, noting that late marriage is framed as harmful for women but not for men, or that ``strong'' personalities are praised in men's charts but moralized in women's. The side-by-side comparison step was the most common ``aha'' moment, because juxtaposing traditional, modern, and reframed columns made gender slippage visible in one glance. At the same time, interviews echoed the Likert pattern on agency: one participant still treated gendered differences as neutral description rather than bias, and several leaned on the modern or AI column as the voice that ``questioned'' the original---raising the design question of whether reframing practice should be required \emph{before} any AI text appears.

The most frequent user-selected bias themes in the quantitative sample (virtuous-wife, chastity, ke-fu, moral judgment) align with these interview accounts and with feminist readings of divination as repetitive gender scripting \parencite{butler1990,foucault1972}. When open reflections defer judgment to ``the AI analyst'' or normalize biased language (``Isn't everyone like this?''), the gap between survey scores and Court text suggests that self-report and performance-based critique may diverge---a pattern consistent with warnings that LLM outputs can function as narrative authorities when framed as analysts rather than critique scaffolds \parencite{bender2021,cotter2022}.

\subsection{Relation to Prior Work}

Workplace bias detectors and engagement-oriented fortune apps leave young readers without a structured reframing practice. \textit{Reframe Destiny} occupies a middle position identified in the literature review: critical without being dismissive, interactive without claiming metaphysical authority. By embedding tagging and reframing inside the same media genre where youth already encounter ke-fu, late-marriage, and ``too strong'' labels, the study tests whether narrative literacy can be taught \emph{in situ} rather than only discussed in the abstract.

\subsection{Limitations}

Limitations include a convenience sample ($N = 19$ paired sessions from 34 total sessions and 23 journey completions), post-survey attrition among journey completers (four missing post-tests), self-report measures, single-session exposure, template-dependent readings, and descriptive rather than inferential analysis. Several sessions show ceiling responding on pre-test items. Interview excerpts are illustrative, not generalizable. Results should not be generalized beyond this sample.

\subsection{Implications and Next Steps}

For design, bias tagging, side-by-side reframing, and a one-sentence rewrite task can be embedded inside divination media in under 15 minutes; interviews suggest adding explicit onboarding for each step. For research, future iterations should require a user-authored Court of Destiny line before showing AI reframe, and should test whether that sequence raises personal agency without reducing awareness gains.

\printbibliography

\end{document}
